\chapter{Conclusão}
\label{ch:conclusao}

\section{Síntese e Contribuições}
\label{sec:conclusao-sintese}

A medicina baseada em dados e a Inteligência Artificial representam a fronteira da inovação em saúde, mas sua viabilidade depende intrinsecamente da superação do dilema central: o equilíbrio entre proteção de privacidade (exigida pela legislação como a LGPD) e preservação de utilidade clínica (essencial para o treinamento de modelos preditivos de alta performance). O presente trabalho investigou este dilema de forma teórica, metodológica e empírica.

\subsection{Consolidação da Fundamentação Teórica e Legal}

A revisão bibliográfica permitiu consolidar uma compreensão profunda da evolução do conceito de privacidade, desde as primeiras formulações de "direito ao isolamento" (Warren e Brandeis, 1890) até as noções contemporâneas de "autodeterminação informativa" (Westin) e "integridade contextual" (Nissenbaum). Ficou evidente que as legislações modernas de proteção de dados — exemplificadas pela LGPD no Brasil — não visam impedir o uso de dados para pesquisa e inovação, mas sim regular seu tratamento de forma responsável.

A pseudonimização robusta emerge como o instrumento técnico privilegiado para viabilizar a pesquisa clínica sob conformidade com a LGPD: ao reduzir o risco de reidentificação a níveis onde a associação a indivíduos específicos não é razoavelmente viável, estas técnicas aproximam os dados do limiar de anonimização efetiva previsto pelo Art. 12 da LGPD, tornando-os tratáveis sob o regime de pesquisa do Art. 11, II. A ANPD, ao consolidar uma abordagem baseada em gestão de risco em seus estudos técnicos sobre anonimização \cite{ANPD2023juridica, ANPD2023computacional}, corrobora que a exigência legal não requer probabilidade zero de reidentificação, mas sua inviabilidade razoável no contexto do tratamento — precisamente o que os experimentos aqui conduzidos demonstram de forma quantitativa. A análise comparativa dos paradigmas de pseudonimização e anonimização (k-anonimato sintático, l-diversidade semântica, e Privacidade Diferencial probabilística) revelou que não existe uma solução universal. Cada técnica oferece garantias distintas com trade-offs próprios: métodos sintáticos como k-anonimato são computacionalmente eficientes mas vulneráveis a ataques de conhecimento; Privacidade Diferencial oferece garantia matemática rigorosa sobre o \textit{release} dos dados mas frequentemente incorre em perda amostral significativa.

\subsection{Desenho e Execução de Metodologia Rigorosa}

O trabalho estabeleceu um delineamento experimental robusto centrado em uma abordagem orientada à tarefa (\textit{task-based evaluation}), onde a utilidade dos dados é medida não por métricas estatísticas genéricas, mas pela preservação da capacidade do modelo preditivo de resolver um problema clínico real. Esta abordagem contrasta com a literatura tradicional, que frequentemente descola a avaliação de utilidade da aplicabilidade prática.

Foi realizado:

\begin{enumerate}
    \item \textbf{Preparação e Caracterização de Dados Reais:} Seleção de uma amostra grande e complexa (42.277 admissões escaladas para 17.917 após filtros de duração) a partir do MIMIC-III, com caracterização demográfica, clínica e distribuição de outcome.
    
    \item \textbf{Definição Operacional de Componentes de Privacidade:} Identificação de quasi-identificadores (Age, Gender, Ethnicity), atributo sensível (Primary Diagnosis), e outcome (Mortality), conforme taxonomia de El Emam.
    
    \item \textbf{Implementação de Três Cenários Progressivos:} Configuração de Scenario A (k-anonymity, $k=5$), Scenario B (k-anonymity + l-diversity, $k=10, l=2$), e Scenario C (Differential Privacy, $\epsilon=2.0$), com especificação clara de hierarquias de generalização e justificativas para cada parâmetro.
    
    \item \textbf{Execução Controlada de Anonimização:} Utilização de ferramenta especializada (ARX) para aplicação determinística dos algoritmos, com registro de métricas de risco (Prosecutor, Journalist, Marketer risks) em cada cenário.
    
    \item \textbf{Protocolo de Avaliação Bifacetado:} Combinação de auditoria de risco de reidentificação (via ARX, três modelos de ameaça) com avaliação de utilidade clínica orientada à tarefa — treinamento dos modelos ConCare e LSTM Late Fusion em dados anonimizados e comparação com baseline nos dados originais, incluindo análise de subgrupos por gênero.
    
    \item \textbf{Rigor Estatístico:} Cálculo de intervalos de confiança de 95\% para métricas de performance (AUROC, AUPRC) através de bootstrap não-paramétrico com 1.000 reamostragens, permitindo quantificação de incerteza.
\end{enumerate}

\subsection{Viabilidade Técnica Consolidada}

A execução completa do pipeline — desde preparação de dados até avaliação de modelos — demonstrou que a aplicação prática de técnicas de anonimização em dados clínicos complexos é viável do ponto de vista técnico. A inclusão de dois modelos com estratégias distintas de integração demográfica — ConCare (\textit{cross-attention}) e LSTM Late Fusion (concatenação tardia) — estabeleceu um teste comparativo rigoroso: o ConCare, por incorporar o vetor demográfico como \textit{query} de atenção sobre toda a série temporal, configura o cenário de maior sensibilidade esperada à anonimização, enquanto o LSTM oferece o contraponto de uma arquitetura com fusão tardia. A convergência na ordenação de robustez entre cenários observada em ambas as arquiteturas fortalece a validade externa das conclusões.

\section{Limitações do Desenho Experimental}
\label{sec:conclusao-limitacoes}

O estudo apresenta limitações que devem ser consideradas na generalização dos achados:

\begin{enumerate}
    \item \textbf{Monocentricidade de Dados e Contexto Demográfico:} A amostra provém exclusivamente do MIMIC-III, coletada de um único hospital (Beth Israel Deaconess Medical Center) nos EUA. Além da monocentricidade geográfica e institucional, o MIMIC-III codifica etnia segundo a taxonomia norte-americana (e.g., "White - Russian", "Black/African American"), que não corresponde ao esquema racial do IBGE utilizado em contextos brasileiros (pardo, preto, branco, amarelo, indígena). As conclusões sobre equidade algorítmica e os parâmetros de hierarquias de generalização para o atributo Ethnicity, portanto, não se transferem diretamente para dados clínicos brasileiros sem adaptação. Generalizações para o SUS ou outros contextos nacionais requerem validação com dados e taxonomias demográficas apropriadas.
    
    \item \textbf{Especificidade Arquitetural:} A avaliação de utilidade clínica foi conduzida sobre dois modelos — ConCare (\textit{cross-attention}) e LSTM Late Fusion (concatenação tardia) — representando os dois paradigmas dominantes de integração de features estáticas e dinâmicas em registros de saúde \cite{Shickel2018}. Embora essa escolha bidirecional aumente a robustez das conclusões em relação a estudos de modelo único, outras arquiteturas (Transformers, GRU-based, modelos sem componente demográfico explícito) podem apresentar sensibilidades distintas à anonimização, e a generalização dos achados para esse espectro requer validação adicional.
    
    \item \textbf{Tarefa Única de Predição:} Testou-se especificamente a tarefa de predição de mortalidade hospitalar. Outras tarefas clínicas (previsão de tempo de permanência, detecção precoce de sepse, recomendação de intervenções) podem apresentar diferentes trade-offs entre privacidade e utilidade.
    
    \item \textbf{Equidade Algorítmica — Análise de Gênero e Lacuna de Etnia:} A análise de subgrupos conduzida estratificou o desempenho por gênero (feminino e masculino), revelando disparidades preexistentes no baseline e efeitos assimétricos das técnicas de pseudonimização entre os subgrupos. No entanto, a supressão de Ethnicity nos Cenários A e B impossibilitou a análise estratificada por raça/etnia — precisamente o eixo demográfico com maior potencial de amplificação de disparidades em algoritmos clínicos \cite{Obermeyer2019}. A avaliação de equidade étnica requer configurações de anonimização que preservem esta variável (como o Cenário C), restringindo as conclusões sobre fairness étnica ao único cenário em que Ethnicity não foi suprimida.
    
    \item \textbf{Parâmetros de Privacidade Diferencial e Escopo da Garantia:} A seleção de $\epsilon=2{,}0$ para Differential Privacy reflete um ponto de equilíbrio empírico reportado na literatura \cite{Pilgram2024}, mas a sensibilidade a variações em $\epsilon$ e $\delta$ não foi sistematicamente explorada. Adicionalmente, a garantia DP aplica-se ao \textit{release} tabular de dados (via ARX/SafePub), não ao modelo treinado; a privacidade diferencial do modelo exigiria mecanismos de treinamento como DP-SGD, fora do escopo deste trabalho.

    \item \textbf{Execuções de Treinamento, Divisão de Dados e Inferência Estatística:} Cada configuração experimental (modelo $\times$ cenário) foi treinada com uma única semente aleatória e avaliada em uma única divisão estática de treino/validação/teste. Diferenças de AUROC de magnitude inferior a 0,02 pontos — como as observadas entre o Baseline e os Cenários B e C — podem estar dentro da variância combinada de treinamento estocástico e partição de dados. A comparação rigorosa entre cenários requereria múltiplas execuções com sementes distintas, validação cruzada estratificada (por subject\_id) e testes pareados (e.g., \textit{bootstrap} pareado sobre $\Delta$AUROC ou teste DeLong), procedimentos não realizados neste trabalho. As conclusões de ordenação entre cenários devem ser interpretadas como indicativas, não como diferenças estatisticamente confirmadas, especialmente nos menores efeitos observados.

    \item \textbf{Escopo da Auditoria de Risco — Quasi-Identificadores Declarados:} A auditoria de risco de reidentificação abrangeu exclusivamente o conjunto de quasi-identificadores declarados \{Age, Gender, Ethnicity\}. O vetor estático dos modelos inclui adicionalmente 25 fenótipos CCS derivados de códigos ICD-9, que podem servir como vetores auxiliares de reidentificação em ataques de ligação com bancos de dados clínicos ou administrativos externos. Combinações raras de comorbidades são conhecidas por seu potencial identificador \cite{Sweeney2002}. Os valores de risco reportados representam, portanto, o risco residual no subconjunto de atributos tratados pela pseudonimização, podendo subestimar o risco real em cenários de ataque com informação clínica auxiliar. Adicionalmente, existe uma tensão conceitual na operacionalização da $l$-diversidade no Cenário B: o atributo sensível protegido (Diagnóstico Primário, derivado de ICD-9) e os fenótipos CCS utilizados como \textit{features} pelo modelo compartilham a mesma fonte de dados (histórico ICD-9), o que significa que o modelo aprende uma representação proxy do atributo que a $l$-diversidade visa proteger. A proteção de $l$-diversidade incide sobre o \textit{dataset} liberado; a exposição indireta via \textit{features} do modelo é uma limitação ortogonal relacionada ao design do vetor de entrada. Por fim, o MIMIC-III foi submetido previamente a de-identificação sob o padrão \textit{Safe Harbor} da HIPAA, de modo que a linha de base de risco 100\% expressa unicidade intra-coorte nos QIs, não exposição a adversários externos com acesso à população geral.

    \item \textbf{Métricas de Calibração:} A avaliação de utilidade clínica concentrou-se nas métricas de discriminação AUROC e AUPRC. Para sistemas de suporte à decisão clínica, a calibração do modelo — medida por Brier Score ou curvas de confiabilidade — é igualmente relevante, pois deslocamentos de distribuição induzidos pela pseudonimização podem degradar a calibração sem impacto proporcional na discriminação. Esta dimensão da utilidade não foi avaliada e constitui direção para trabalhos futuros.
\end{enumerate}

\section{Considerações Finais}
\label{sec:conclusao-finais}

Este trabalho responde de forma operacional e rigorosa à pergunta de pesquisa central: \textbf{A anonimização pode viabilizar o uso legal de dados sensíveis sob a LGPD, preservando a utilidade clínica de modelos de Deep Learning em UTI?} Através de uma abordagem teórico-metodológica-empírica integrada, demonstra-se que a resposta não é binária. Não existe uma única técnica que simultaneamente maximiza privacidade e utilidade; ao contrário, há um espectro de escolhas com trade-offs explícitos.

O trabalho oferece diretrizes técnicas concretas:

\begin{itemize}
    \item A aplicação de técnicas de anonimização robustas (Cenário B: $k=10, l=2$) reduz o risco de reidentificação para níveis estatisticamente negligenciáveis — compatíveis com o padrão de ``meios técnicos razoáveis'' do Art. 12 da LGPD \cite{Brasil2018lgpd} e com a interpretação regulatória da ANPD \cite{ANPD2023juridica} —, mantendo viabilidade para modelos preditivos complexos;
    \item A preservação de tamanho amostral é crítica: neste estudo, a Privacidade Diferencial suprimiu 31\% dos registros de treinamento, o que confunde seu impacto na utilidade com o efeito de simplesmente treinar com menos dados. Técnicas sintáticas robustas que mantêm o volume amostral intacto (Cenário B: $k=10, l=2$) apresentaram melhor utilidade do que a Privacidade Diferencial com supressão amostral expressiva; a generalização deste achado para configurações de DP sem supressão requer investigação adicional;
    \item A abordagem de avaliar utilidade através de tarefa clínica específica (não apenas métricas genéricas) fornece uma noção mais rica e aplicável de "utilidade aceitável".
\end{itemize}

Estes resultados constituem evidência técnica quantitativa para informar gestores de saúde, pesquisadores e formuladores de políticas públicas sobre o uso responsável de dados sensíveis, contribuindo para a construção de um arcabouço técnico-regulatório alinhado com a LGPD e com o avanço da inovação em saúde.

\section{Trabalhos Futuros}
\label{sec:conclusao-futuros}

As seguintes linhas de investigação emergem como naturais prolongamentos deste trabalho:

\begin{enumerate}
    \item \textbf{Validação em Dados Brasileiros:} Replicação do desenho experimental em bases de dados clínicos brasileiros (ex: dados do SUS, sistemas de prontuário eletrônico de hospitais brasileiros) para validar se os trade-offs privacidade-utilidade permanecem semelhantes em contextos nacionais com diferentes distribuições demográficas e epidemiológicas.
    
    \item \textbf{Análise de Equidade Algorítmica:} Investigação sistemática de como a anonimização (especialmente supressão de Ethnicity) impacta a distribuição de erros entre grupos demográficos, utilizando métricas de fairness (disparate impact, equalized odds).
    
    \item \textbf{Generalização Arquitetural e de Técnicas de Privacidade:} Teste do impacto da anonimização em modelos alternativos (Transformers, GRU-based architectures, modelos sem componente demográfico explícito) para determinar se as conclusões transcendem as duas arquiteturas avaliadas. Paralelamente, a inclusão de $t$-closeness como quarto cenário experimental permitiria completar o espectro de técnicas sintáticas discutidas na Fundamentação Teórica — atualmente a técnica é descrita teoricamente mas não avaliada empiricamente.
    
    \item \textbf{Otimização de Parâmetros:} Exploração sistemática da sensibilidade dos trade-offs a variações em $k$, $l$, $\epsilon$ e configurações de hierarquias de generalização, potencialmente através de técnicas de busca automática de hiperparâmetros.
    
    \item \textbf{Integração Jurídica e Regulatória:} Diálogo aprofundado com especialistas em direito da proteção de dados para determinar qual limiar operacional de risco é aceito pelo texto e pela jurisprudência da LGPD, transformando estes achados técnicos em recomendações normativas concretas. Os estudos técnicos da ANPD sobre anonimização \cite{ANPD2023juridica, ANPD2023computacional} constituem o ponto de partida regulatório mais atualizado, mas carecem de vinculação com limiares numéricos específicos de risco residual — lacuna que este trabalho começa a preencher empiricamente e que futuros estudos podem endereçar com maior rigor jurídico-normativo.

    \item \textbf{Avaliação de Calibração e Testes Pareados:} Extensão da avaliação de utilidade para incluir métricas de calibração (Brier Score, curvas de confiabilidade) e realização de testes estatísticos pareados entre cenários ($\Delta$AUROC via bootstrap pareado ou teste DeLong) para confirmar com rigor as diferenças de desempenho observadas entre cenários.

    \item \textbf{Privacidade Diferencial no Treinamento:} Investigação do impacto de mecanismos de treinamento com DP (DP-SGD) aplicados diretamente ao processo de otimização dos modelos ConCare e LSTM, comparando com a abordagem de DP sobre o \textit{release} de dados adotada neste trabalho.

    \item \textbf{Avaliação de Ataques de Inferência de Pertencimento (\textit{Membership Inference}):} Este trabalho avaliou o risco de reidentificação nos dados anonimizados, mas não a privacidade dos \textit{modelos treinados}. Ataques de inferência de pertencimento (\textit{membership inference attacks}) — que determinam se um registro específico foi utilizado no treinamento do modelo — representam uma ameaça de privacidade distinta e crescente em modelos de \textit{Deep Learning} aplicados a dados clínicos \cite{Olatunji2024}. Investigar a vulnerabilidade dos modelos ConCare e LSTM a este tipo de ataque, e se a anonimização dos dados de treinamento reduz essa vulnerabilidade, constitui uma extensão natural deste trabalho.
\end{enumerate}

Espera-se que este trabalho contribua não apenas para o avanço da ciência da computação e da informática em saúde, mas também para a viabilização prática de pesquisa clínica sob conformidade legal no Brasil.